\documentclass[conference]{IEEEtran}
\IEEEoverridecommandlockouts
% The preceding line is only needed to identify funding in the first footnote. If that is unneeded, please comment it out.
\usepackage{cite}
\usepackage{amsmath,amssymb,amsfonts}
\usepackage{algorithmic}
\usepackage[T1]{fontenc}
\usepackage[utf8]{inputenc}
\usepackage{tabularx}
\usepackage{booktabs}
\usepackage{url}
\usepackage{graphicx}
\usepackage{listings}
\usepackage{listingsutf8}
\usepackage{pifont}
\usepackage{textcomp}
\usepackage{xcolor}
\usepackage{hyperref}
\usepackage[german]{babel}
\usepackage{blindtext}
\usepackage{tikz}
\usepackage{lstautogobble}
\usepackage[framemethod=TikZ]{mdframed}
\usepackage{bookmark}

\usetikzlibrary{calc,arrows.meta,decorations,backgrounds, fit}
\def\BibTeX{{\rm B\kern-.05em{\sc i\kern-.025em b}\kern-.08em
    T\kern-.1667em\lower.7ex\hbox{E}\kern-.125emX}}

\lstset{
  basicstyle=\small\ttfamily,
  breaklines=true,
  frame=single,
  backgroundcolor=\color{gray!20},
  literate=
    {→}{{$\rightarrow$}}1
    {←}{{$\leftarrow$}}1
    {⇒}{{$\Rightarrow$}}1
    {⇐}{{$\Leftarrow$}}1
}

\lstdefinestyle{consolestyle}{
  basicstyle=\color{green!90!black!100}\ttfamily\footnotesize,
  backgroundcolor=\color{black!80},
  frame=single,
  breaklines=true,
  columns=fullflexible,
  keepspaces=true,
  autogobble=true
}

\newcommand{\cmark}{\ding{51}}
\newcommand{\xmark}{\ding{55}}
\newcommand{\warnmark}{$\circ$}

\begin{document}

\title{TLS - Transport Layer Security\\
{\footnotesize Eine Betrachtung von TLS hinsichtlich Sicherheit und Robustheit }
}

\author{\IEEEauthorblockN{1\textsuperscript{st} Hannah Herberstein}
\IEEEauthorblockA{\textit{HKA} \\
\textit{Hochschule Karlsruhe}\\
heha1015@h-ka.de}
\and
\IEEEauthorblockN{2\textsuperscript{nd} Leonie Leber}
\IEEEauthorblockA{\textit{HKA} \\
\textit{Hochschule Karlsruhe}\\
lele1018@h-ka.de}
\and
\IEEEauthorblockN{3\textsuperscript{rd} Sina Frieß}
\IEEEauthorblockA{\textit{HKA} \\
\textit{Hochschule Karlsruhe}\\
frsi1018@h-ka.de}
\and
\IEEEauthorblockN{4\textsuperscript{th} Niels Reisch}
\IEEEauthorblockA{\textit{HKA} \\
\textit{Hochschule Karlsruhe}\\
reni1030@h-ka.de}
}

\maketitle

\begin{abstract}

Transport Layer Security (TLS) bildet das kryptographische Fundament sicherer Netzwerkkommunikation. Diese Arbeit analysiert den TLS-Handshake-Prozess, die Zertifikatsvalidierung sowie Vertrauensketten anhand zweier kontrastierender Client-Server-Szenarien: einer gehärteten TLS~1.3-Konfiguration mit Perfect Forward Secrecy und einem System mit veralteter Protokollversion und schwachen Cipher Suites. Am Beispiel des POODLE-Angriffs (CVE-2014-3566) wird ein gezielter Protokoll-Downgrade auf SSL~3.0 simuliert und anhand synthetischer Handshake-Logs nachvollzogen, an welcher Stelle die Sicherheitsmechanismen versagen. Ergänzend wird das Szenario eines selbstsignierten Zertifikats untersucht, das den Handshake mit einem \texttt{CERTIFICATE\_UNKNOWN}-Alert abbricht und die Verwundbarkeit gegenüber Impersonation-Angriffen aufzeigt. Die Bewertung der Auswirkungen auf Vertraulichkeit, Integrität und Authentizität zeigt, dass Legacy-Protokolle alle drei Schutzziele gleichermaßen gefährden. Als Härtungsmaßnahmen werden die ausschließliche Verwendung von TLS~1.2 und TLS~1.3, der Einsatz von Cipher Suites mit Perfect Forward Secrecy, OCSP Stapling, Certificate Transparency sowie HTTP Strict Transport Security diskutiert und im Einklang mit der BSI-Technischen Richtlinie TR-02102-2 bewertet.

\end{abstract}

\begin{IEEEkeywords}
Transport Layer Security, TLS Handshake, Downgrade-Angriff,
POODLE, Cipher Suite, Perfect Forward Secrecy, Zertifikatsvalidierung,
Public Key Infrastructure, Härtungsmaßnahmen, IT-Sicherheit
\end{IEEEkeywords}

\section{Einführung}
\subsection{Überblick über TLS}
Transport Layer Security (TLS) ist ein kryptografisches Sicherheitsprotokoll zur Absicherung der Datenübertragung in Netzwerken \\cite{rfc8446}. Es wird eingesetzt, um die Kommunikation zwischen zwei Parteien, Client und Server, vor Abhören, Manipulation und Identitätsfälschung zu schützen. TLS ist heute der Standard zur Absicherung von TCP-basierten Anwendungen und wird unter anderem im Internetprotokoll HTTPS verwendet. TLS ist der Nachfolger von Secure Socket Layer (SSL). Obwohl häufig noch von „SSL“ gesprochen wird, kommen in modernen Systemen ausschließlich TLS-Versionen zum Einsatz. TLS arbeitet zwischen der Anwendungs- und der Transportschicht und ermöglicht eine sichere Kommunikation über unsichere Netzwerke wie das Internet.
\subsection{Ziele von TLS}
TLS verfolgt drei wesentliche Sicherheitsziele:
\begin{itemize}
  \item \textbf{Vertraulichkeit (Confidentiality):} Schutz der übertragenen Daten vor unbefugtem Mitlesen durch Verschlüsselung.
  \item \textbf{Integrität (Integrity):} Sicherstellung, dass Daten während der Übertragung nicht verändert oder manipuliert wurden.
  \item \textbf{Authentifizierung (Authentication):} Überprüfung der Identität des Kommunikationspartners mithilfe digitaler Zertifikate.
\end{itemize}

\subsection{Grundlegende Funktionsweise von TLS}
TLS kombiniert verschiedene kryptografische Verfahren:

\subsection*{Symmetrische Verschlüsselung}
Die eigentliche Datenübertragung erfolgt mithilfe symmetrischer Verschlüsselung. Dabei verwenden Client und Server denselben geheimen Sitzungsschlüssel (Session Key). Symmetrische Verfahren sind sehr effizient und eignen sich daher besonders für die schnelle Verschlüsselung großer Datenmengen.

\subsection*{Asymmetrische Verschlüsselung}
Zur sicheren Aushandlung bzw. zum Schutz der Sitzungsschlüssel werden asymmetrische Verfahren eingesetzt. Hierbei existiert ein Schlüsselpaar aus Public Key und Private Key.
In modernen TLS-Versionen erfolgt der Schlüsselaustausch hauptsächlich über ECDHE (Elliptic Curve Diffie-Hellman Ephemeral) \\cite{rfc8446}. Dadurch kann ein gemeinsamer geheimer Schlüssel erzeugt werden, ohne dass dieser direkt übertragen werden muss.

\subsection*{Integritätsschutz}
TLS stellt sicher, dass Daten während der Übertragung nicht manipuliert werden. Ältere TLS-Versionen verwendeten dafür MACs (Message Authentication Codes). Moderne TLS-Versionen, insbesondere TLS 1.3, nutzen überwiegend sogenannte AEAD-Verfahren (Authenticated Encryption with Associated Data), die Verschlüsselung und Integritätsschutz kombinieren \\cite{rfc8446}.

\subsection*{Authentifizierung}
Die Authentifizierung erfolgt über digitale Zertifikate. Diese bestätigen die Identität des Servers und ermöglichen dem Client die Überprüfung, ob er tatsächlich mit dem gewünschten Kommunikationspartner verbunden ist.


\subsection{Die TLS-Verbindung}
\subsection*{Verbindungsaufbau}
Eine TLS-Verbindung besteht im Wesentlichen aus vier Phasen:
\begin{enumerate}
  \item TLS-Handshake
  \item Authentifizierung und Zertifikatsprüfung
  \item Schlüsselaustausch und Schlüsselableitung
  \item Verschlüsselte Datenübertragung
\end{enumerate}

In der Abbildung~\ref{fig:tls13-handshake} im Anhang (Abschnitt~\ref{app:tls-handshake}) zum TLS-1.3-Handshake ist der grundlegende Ablauf des Verbindungsaufbaus zwischen Client und Server dargestellt. Zunächst sendet der Client eine Client Hello-Nachricht an den Server. Diese enthält unter anderem die unterstützte TLS-Version, eine Liste unterstützter Cipher Suites, Zufallswerte (Client Random) sowie Parameter für den Schlüsselaustausch. \\
Der Server antwortet anschließend mit einer Server Hello-Nachricht. Darin teilt der Server die ausgewählte Cipher Suite mit und übermittelt sein digitales Zertifikat sowie eine digitale Signatur. Mithilfe des Zertifikats kann sich der Server gegenüber dem Client authentifizieren. Der Client überprüft daraufhin die Gültigkeit des Zertifikats sowie die Signatur der Zertifizierungsstelle. \\
Nach erfolgreicher Verifikation führen beide Kommunikationspartner den Schlüsselaustausch durch und erzeugen daraus einen gemeinsamen Sitzungsschlüssel (Session Key). Dieser Schlüssel wird anschließend für die symmetrische Verschlüsselung der weiteren Kommunikation verwendet. Nachdem der Handshake abgeschlossen wurde, kann die eigentliche verschlüsselte Datenübertragung beginnen.


\subsection*{Cipher Suites}
Eine Cipher Suite beschreibt die kryptografischen Verfahren, die während der TLS-Verbindung verwendet werden. Dazu gehören unter anderem:
\begin{itemize}
  \item Verfahren zum Schlüsselaustausch
  \item Verschlüsselungsalgorithmen
  \item Verfahren zur Integritätsprüfung
\end{itemize}
Client und Server müssen sich während des Handshakes auf eine gemeinsame Cipher Suite einigen.

\subsubsection*{Unterschiede zu älteren TLS-Versionen}
In älteren TLS-Versionen wurde häufig RSA für den Schlüsselaustausch verwendet. Dabei erzeugte der Client einen geheimen Sitzungsschlüssel und verschlüsselte diesen mit dem öffentlichen Schlüssel des Servers. \\
TLS 1.3 unterstützt dieses Verfahren nicht mehr. Stattdessen wird überwiegend ECDHE verwendet, da dieses Verfahren sicherer ist und sogenannte Perfect Forward Secrecy ermöglicht. Dadurch bleiben vergangene Kommunikationsdaten selbst dann geschützt, wenn ein langfristiger Schlüssel kompromittiert wird.

\subsection*{Authentifizierung des Servers}
Der Server authentifiziert sich gegenüber dem Client mithilfe eines digitalen Zertifikats. Dieses Zertifikat enthält unter anderem:
\begin{itemize}
  \item den Domainnamen
  \item den öffentlichen Schlüssel des Servers
  \item der Gültigkeit
  \item Informationen über die ausstellende Zertifizierungsstelle
  \item die digitale Signatur der Certificate Authority (CA)
\end{itemize}
Während des TLS-Handshakes weist der Server nach, dass er im Besitz des zugehörigen Private Keys ist. Der Client überprüft anschließend die digitale Signatur des Zertifikats und kontrolliert, ob dieses von einer vertrauenswürdigen Zertifizierungsstelle ausgestellt wurde. Dadurch kann der Client sicherstellen, dass der Server authentisch ist.

\subsection*{Zertifikatsprüfung und Vertrauensketten}
Die grundlegende Idee eines TLS-Zertifikats besteht darin, einen öffentlichen Schlüssel eindeutig mit einer Identität - meist einem Domainnamen - zu verknüpfen. Der Client vertraut dabei nicht direkt dem Serverzertifikat selbst, sondern der Zertifizierungsstelle, die dieses Zertifikat signiert hat.

Die Vertrauenskette bei TLS sieht typischerweise wie folgt aus:
\begin{lstlisting}
Root CA → signiert → Intermediate CA → signiert → Serverzertifikat
\end{lstlisting}
Die Root Certificate Authority ist selbstsigniert und befindet sich bereits im sogenannten Trust Store des Betriebssystems oder Browsers. Dieser Trust Store enthält vertrauenswürdige Root-Zertifikate. \\
Aus Sicherheitsgründen signieren Root CAs Serverzertifikate normalerweise nicht direkt. Stattdessen werden sogenannte Intermediate CAs verwendet. Dies bietet mehrere Vorteile:
\begin{itemize}
  \item höhere Sicherheit
  \item einfacherer Widerruf kompromittierter Zertifikate
  \item geringeres Risiko bei einer Kompromittierung der Root CA
\end{itemize}
Die Intermediate CA signiert anschließend das eigentliche Serverzertifikat.

\subsubsection*{Überprüfung der Zertifikatskette}
Beim Verbindungsaufbau überprüft der Browser bzw. Client mehrere Punkte:
\begin{itemize}
  \item Ist die digitale Signatur korrekt?
  \item Ist die Zertifikatskette vollständig?
  \item Ist die Root CA vertrauenswürdig?
  \item Ist das Zertifikat noch gültig?
  \item Passt der Domainname zum Zertifikat?
\end{itemize}
Kann der Client die Zertifikatskette erfolgreich bis zu einer vertrauenswürdigen Root CA zurückverfolgen, wird das Zertifikat als vertrauenswürdig akzeptiert und die Verbindung hergestellt.

\subsection*{Schlüsselaustausch und Session Keys}
Nach erfolgreicher Authentifizierung führen Client und Server einen kryptografischen Schlüsselaustausch durch. In modernen TLS-Versionen erfolgt dies meist über ECDHE. Dabei erzeugen beide Kommunikationspartner gemeinsam einen geheimen Sitzungsschlüssel, ohne diesen direkt zu übertragen. Aus diesem Schlüsselaustausch werden anschließend die Session Keys für die symmetrische Verschlüsselung abgeleitet.

\subsection*{Verschlüsselte Datenübertragung}
Nachdem der Handshake abgeschlossen wurde und beide Seiten gemeinsame Sitzungsschlüssel besitzen, beginnt die eigentliche Kommunikation. Die übertragenen Daten werden nun symmetrisch verschlüsselt. Zusätzlich wird ihre Integrität überprüft.

\subsection*{Integritätsschutz mit MACs und AEAD}
Bei älteren TLS-Versionen wurde ein Message Authentication Code (MAC) verwendet. Dabei berechnen Sender und Empfänger mithilfe des gemeinsamen geheimen Schlüssels einen Prüfwert über die übertragenen Daten. Stimmen die berechneten Werte überein, wurde die Nachricht nicht manipuliert. Moderne TLS-Versionen verwenden überwiegend AEAD-Verfahren wie AES-GCM oder ChaCha20-Poly1305. Diese kombinieren Verschlüsselung und Integritätsschutz in einem einzigen Verfahren.

\subsection{Zusammenfassung}
TLS ermöglicht eine sichere Kommunikation über unsichere Netzwerke. Dazu kombiniert das Protokoll asymmetrische und symmetrische Kryptografie, digitale Zertifikate sowie kryptografische Integritätsprüfungen.
Der TLS-Handshake dient dabei der sicheren Aushandlung aller benötigten Parameter und Schlüssel. Zertifikate und Vertrauensketten ermöglichen die Authentifizierung des Servers. Anschließend erfolgt die verschlüsselte und integritätsgeschützte Datenübertragung mithilfe gemeinsamer Sitzungsschlüssel.

\section{Konfigurationen von TLS}
Dieses Kapitel beschreibt ein Testszenario, das ein gehärtetes und ein verwundbares TLS-System gegenüberstellt. 

\subsection{Sichere Konfiguration}
Client-Server-Paar~1 simuliert eine moderne, sichere TLS-Konfiguration. Die wichtigsten Parameter sind Protokollversion, Schlüsselaustausch, Verschlüsselungsmodus, Integritätsschutz, Zertifikatstyp und Schlüssellänge.

Als Protokollversion wird TLS~1.3 verwendet~\cite{rfc8446}. Der Schlüsselaustausch erfolgt über ECDHE, sodass Perfect Forward Secrecy erreicht wird. Verschlüsselung und Integritätsschutz setzt das Paar über AEAD im GCM-Modus um~\cite{gcm-wiki}. Als Zertifikat dient ein CA-signiertes Zertifikat mit ECDSA oder 4096-Bit RSA~\cite{ecdsa-ref,bsi-tr02102}.

\subsection{Veraltete Konfiguration}
Das Client-Server-Paar~2 stellt ein ungepatchtes Altsystem dar und dient als Primärziel für die späteren Fehler-Simulationen. Es beinhaltet die unsichere TLS-Konfiguration.

Es nutzt TLS~1.2 mit rückwärtskompatiblen, unsicheren Mechanismen~\cite{rfc8996}. Der Schlüsselaustausch erfolgt statisch per RSA~\cite{rsa-studyflix}, wodurch keine Forward Secrecy besteht. Verschlüsselung per CBC und Integritätsschutz per SHA-1 MAC sind anfällig für Padding-Angriffe. Das Zertifikat ist selbstsigniert und verwendet nur 1024-Bit RSA~\cite{bsi-tr02102}.

\subsection{Konfiguration der Simulationen}

Die exakten OpenSSL-Befehle und das Setup sind im Anhang (Abschnitt~\ref{app:simulation-config}) dokumentiert. Dort finden Sie die Zertifikatsgenerierung, Client-Server-Setups sowie die Szenarien für erfolgreichen und nicht-erfolgreichen Downgrade.
\section{Experimentelles Setup und Methodik}

\subsection{Simulationsumgebung}

Für die praktische Analyse wurde eine lokale Loopback-Umgebung (127.0.0.1) auf einem
Windows-System eingerichtet. OpenSSL 3.6.1 diente auf Client und Server als kryptographische Grundlage~\cite{openssl}. Der Datenverkehr wurde parallel mit Wireshark aufgezeichnet, um den Handshake-Verlauf und Fehlermeldungen auf Paketebene zu analysieren.

\subsection{TLS-Fehlersimulation und Log-Analyse}

Es werden ein Downgrade- und ein Handshake-Fehlerfall simuliert, um zu untersuchen, wo Schutzmechanismen greifen oder versagen.

% ── Szenario 1 ────────────────────────────────────────────────────────────────
\subsection{Szenario~1: Nicht-erfolgreicher Downgrade (Handshake-Fehlerfall)}

Im ersten Experiment wurde untersucht, wie ein moderner Client reagiert, wenn ein Server
ein veraltetes, unsicheres Protokoll erzwingt. Dies simuliert das Verhalten bei einer
serverseitigen Fehlkonfiguration oder einem Manipulationsversuch im
Netzpfad~\cite{sentinelone}.

\subsection*{Vorbereitung der Zertifikate}

Zunächst wurde durch die Nutzung von OpenSSL ein selbstsigniertes X.509-Zertifikat mit
einem 2048-Bit RSA-Schlüsselpaar generiert. Der Parameter \texttt{-nodes} (No-DES) sorgt
dafür, dass der private Schlüssel unverschlüsselt bleibt, um einen reibungslosen
Serverstart ohne Passwortabfrage zu ermöglichen~\cite{openssl}:

\begin{lstlisting}[style=consolestyle]
openssl req -x509 -newkey rsa:2048 -keyout key.pem -out cert.pem \
-days 365 -nodes
\end{lstlisting}

\subsection*{Verbindungsaufbau}

Der simulierte Server wurde so gestartet, dass er ausschließlich das veraltete Protokoll
TLS~1.0 anbietet. Anschließend versuchte der reguläre Client, eine sichere Verbindung
zu diesem Port aufzubauen. Siehe Anhang~\ref{app:sz1-befehle}.

\subsection*{Analyse der Handshake-Schritte}

Der Verbindungsaufbau wird durch den Client gestartet, worauf der Server antwortet und
versucht, TLS~1.0 aufzuzwingen~\cite{ssldragon}. An dieser Stelle greift der
clientseitige Sicherheitsmechanismus: Moderne OpenSSL-Bibliotheken blockieren
standardmäßig die Kommunikation über Protokolle unterhalb von TLS~1.2~\cite{openssl}.
Der Handshake bricht unmittelbar ab. Das Terminal des Clients dokumentiert diesen fatalen
Protokollkonflikt wie in Anhang~\ref{app:sz1-log} aufgezeigt.

In der parallelen Wireshark-Aufzeichnung zeigt sich dieser Abbruch durch ein rot
markiertes Paket mit der Kennzeichnung
\texttt{Alert (Level: Fatal, Description: Protocol Version (70))}.
Der Fehlercode~70 (\textit{protocol\_version}) belegt, dass der Client die erzwungene
Abwärtskompatibilität des Servers zum Schutz des Anwenders blockiert
hat~\cite{alertcode70}.

% ── Szenario 2 ────────────────────────────────────────────────────────────────
\subsection{Szenario~2: Erfolgreicher Downgrade (Schwächung der Mechanismen)}

Das zweite Experiment simuliert den Fall, in dem die clientseitigen Schutzmechanismen
bewusst deaktiviert werden -- beispielsweise, um die Kommunikation mit ungenügend
gehärteten Altsystemen (Legacy-Systemen) im Firmennetzwerk zu
erzwingen~\cite{sentinelone}.

\subsection*{Verbindungsaufbau unter reduzierten Sicherheitsbarrieren}

Nun werden der Server und der Client durch das explizite Herabsetzen des
Sicherheitslevels auf die Stufe~Null (\texttt{@SECLEVEL=0}) jeweils gezwungen, ihre
Schutzschilde zu deaktivieren~\cite{openssl}~(Anhang~\ref{app:sz2-befehle}).

\subsection*{Analyse der Handshake-Schritte}

Durch die Modifikation signalisiert der Client bereits im \textit{Client Hello}, dass er
im Notfall auch das unsichere TLS~1.0 akzeptiert~\cite{ssldragon}. Der Server antwortet
im \textit{Server Hello} mit der Bestätigung der veralteten Version. Da der Client das
nicht mehr blockiert, wechselt die Verbindung in den Zustand \texttt{CONNECTED}.
Anhang~\ref{app:sz2-log} zeigt das Protokoll-Log des Terminals des Clients.

\subsection*{Bewertung der Schwachstelle}

Mit abgeschwächten Sicherheitsregeln ist die Verbindung zwar aufgebaut (\texttt{CONNECTED}), aber nicht mehr
zeitgemäß geschützt. TLS~1.0 bietet keine ausreichende Vertraulichkeit und Integrität,
so dass ein Man-in-the-Middle-Angriff möglich wird~\cite{sentinelone,mdndocs}.

% ── Vergleichstabelle ─────────────────────────────────────────────────────────
\subsection{Analytische Gegenüberstellung der Simulationsergebnisse}

Zur übersichtlichen Strukturierung der Ergebnisse sind die Auswirkungen der beiden Konfigurationen in Tabelle~\ref{tab:vergleich} im Anhang (Abschnitt~\ref{app:vergleichstabelle}) zusammenfassend gegenübergestellt.

\section{Auswirkungen eines Fehlerfalls auf Schutzziele}

Die Auswirkungen der beiden untersuchten Szenarien sind in Tabelle~\ref{tab:Auswirkungen-Schutzziele} im Anhang (Abschnitt~\ref{app:auswirkungstabelle}) zusammengefasst.

\subsection*{Angriff POODLE} 
Der POODLE-Angriff (Padding Oracle On Downgraded Legacy Encryption, CVE-2014-3566) zeigt, wie ein Angreifer eine Verbindung von TLS auf SSL~3.0 downgraded und dadurch die Verbindung kompromittiert \cite{cve-poodle}. Die Padding-Oracle-Eigenschaft von SSL~3.0 im CBC-Modus ermöglicht die iterative Rekonstruktion von Klartextbytes durch einen Angreifer im Netzpfad (\textit{Vertraulichkeit}). Manipulierte CBC-Blöcke werden akzeptiert, sofern das Padding zufällig korrekt ist -- die MAC-then-Encrypt-Konstruktion bietet keinen ausreichenden Schutz (\textit{Integrität}). Mit dem erzwungenen Wechsel auf SSL~3.0 entfallen die Finished-Message-Signaturen moderner TLS-Versionen; ohne \texttt{TLS\_FALLBACK\_SCSV} bleibt der Downgrade unbemerkt (\textit{Authentizität}) \cite{thomas-krenn-poodle}.

\subsection*{Selbstsigniertes Zertifikat} 

Präsentiert ein Server ein selbstsigniertes Zertifikat, bricht ein korrekt konfigurierter Client den Handshake mit einem \texttt{CERTIFICATE\_UNKNOWN}-Alert ab -- die Verbindung kommt nicht zustande \cite{mdndocs}. Deaktiviert ein Nutzer diese Prüfung jedoch (z.\,B.\ durch Akzeptieren einer Browser-Warnung oder mittels \texttt{-{}-no-verify}-Flag), öffnet sich ein Angriffsfenster: Ein Angreifer im Netzpfad kann ein eigenes selbstsigniertes Zertifikat einschleusen, das der Client mangels CA-Prüfung ebenfalls akzeptiert. Damit terminiert der Angreifer \emph{beide} TLS-Verbindungen separat -- zum Client und zum echten Server -- und leitet Daten transparent weiter (\textit{Man-in-the-Middle}). Alle übertragenen Klartextdaten sind dem Angreifer zugänglich, obwohl der Nutzer eine \texttt{https://}-Verbindung im Browser sieht (\textit{Vertraulichkeit verletzt, Authentizität verletzt}). Die TLS-Record-MAC-Prüfung schützt die Integrität \emph{innerhalb} jeder der beiden Teilverbindungen -- der Angreifer kann jedoch Inhalte modifizieren, bevor er sie in der zweiten Verbindung weitersendet, sodass der Endnutzer manipulierte Daten als integer wahrnimmt (\textit{Integrität de facto verletzt}).

\section{Härtungsmaßnahmen}
\subsection*{TLS Versionsbeschränkung}

Das BSI schreibt in der \textbf{Technischen Richtlinie TR-02102-2} (Stand: Januar 2026) vor, dass in neuen Systemen ausschließlich TLS~1.2 und TLS~1.3 eingesetzt werden sollen \cite{bsi-tr02102}. SSL~2.0, SSL~3.0 und TLS~1.0 sind vollständig zu deaktivieren; TLS~1.1 wird ebenfalls nicht mehr empfohlen. Die Konfiguration in nginx lautet entsprechend:

\begin{lstlisting}
ssl_protocols TLSv1.2 TLSv1.3;
\end{lstlisting}

Durch die Abschaltung aller Legacy-Protokolle entfällt die Angriffsfläche für POODLE und BEAST \cite{cve-beast} grundlegend, da kein Protokoll-Fallback mehr stattfinden kann \cite{thomas-krenn-poodle}. Ergänzend sollte \textbf{TLS\_FALLBACK\_SCSV} (RFC~7507) implementiert werden: Übermittelt ein Client diesen Signaling Cipher Suite Value in einem erneuten, herabgestuften Verbindungsversuch, bricht der Server den Handshake mit einem \texttt{inappropriate\_fallback}-Alert ab und verhindert so auch aktive Downgrade-Versuche bei Systemen, die Legacy-Versionen aus Kompatibilitätsgründen noch führen \cite{rfc7507}.

\subsection*{Sichere Cipher Suites}

Algorithmen ohne kryptographische Langzeitsicherheit -- darunter RC4, DES, 3DES, Export-Cipher und MD5-basierte MACs -- sind vollständig aus der Konfiguration zu entfernen \cite{bsi-tr02102}. Gemäß BSI~TR-02102-2 sind ausschließlich Cipher Suites mit \textbf{Perfect Forward Secrecy (PFS)} zu verwenden: Nur ephemere Schlüsselaustauschverfahren (ECDHE, DHE) stellen sicher, dass ein nachträglich kompromittierter Langzeitschlüssel keine bereits aufgezeichneten Sessions entschlüsselbar macht. Empfohlene Suiten für TLS~1.2 sind:

\begin{lstlisting}
TLS_ECDHE_RSA_WITH_AES_256_GCM_SHA384
TLS_ECDHE_ECDSA_WITH_AES_128_GCM_SHA256
TLS_DHE_RSA_WITH_AES_256_GCM_SHA384
\end{lstlisting}

Für TLS~1.3 sind die verfügbaren Cipher Suites protokollseitig bereits auf AEAD-Verfahren (AES-GCM, ChaCha20-Poly1305) beschränkt; RC4, CBC sowie der statische RSA-Schlüsseltransport sind im Standard nicht mehr vorgesehen \cite{windows-papst-ciphers,mozilla-ssl-gen}.

\subsection*{Zertifikatsvalidierung und Vertrauensketten}

Selbstsignierte Zertifikate sind in produktiven Umgebungen unzulässig. Es sind Zertifikate einer anerkannten Zertifizierungsstelle einzusetzen, deren Root-Zertifikate in den Betriebssystem- und Browser-Trust-Stores enthalten sind. Zur weiteren Absicherung der Vertrauenskette empfehlen sich folgende Mechanismen:

\begin{itemize}
  \item \textbf{Certificate Transparency (CT):} 
    Alle ausgestellten Zertifikate werden in öffentliche, append-only-Logs eingetragen. Clients und Monitore können überprüfen, ob ein präsentiertes Zertifikat dort verzeichnet ist, und unautorisierte Ausstellungen frühzeitig erkennen.
  \item \textbf{OCSP Stapling:} 
    Der Server liefert den aktuellen Sperrstatus seines Zertifikats unmittelbar im TLS-Handshake mit, sodass der Client keinen separaten OCSP-Request senden muss. Dies     verhindert, dass ein Angreifer OCSP-Abfragen blockiert und ein bereits widerrufenes Zertifikat weiterhin nutzbar bleibt.
  \item \textbf{CAA-DNS-Einträge:} 
    \emph{Certification Authority Authorization} Records schränken auf DNS-Ebene ein, welche Zertifizierungsstellen für eine Domain überhaupt Zertifikate ausstellen dürfen, und verringern so das Risiko unberechtigter Ausstellungen.
\end{itemize}

\subsection*{Regelmäßige Überprüfung und Patch-Management}

Da Schwachstellen kontinuierlich neu entdeckt werden (vgl. MITRE CVE-Datenbank), ist eine \textbf{regelmäßige Überprüfung der TLS-Konfiguration} zwingend erforderlich. Werkzeuge wie \href{https://testssl.sh}{testssl.sh} oder der \href{https://www.ssllabs.com/ssltest/}{Qualys SSL~Labs Server Test} erlauben eine automatisierte Analyse aktiver Protokollversionen, Cipher Suites und Zertifikatseigenschaften. Das BSI macht die TR-02102-2 für Bundesbehörden gemäß § 8 Abs.~1 BSIG verbindlich und aktualisiert die Richtlinie fortlaufend -- zuletzt im Januar 2026 \cite{kes-bsi-tr}.


\appendices
\section{OpenSSL-Befehle Szenario~1}
\label{app:sz1-befehle}

\noindent\textbf{Server-Befehl (TLS~1.0 erzwingen):}
\begin{lstlisting}[style=consolestyle]
openssl s_server -accept 4433 -cert cert.pem -key key.pem -tls1
\end{lstlisting}

\noindent\textbf{Client-Befehl (Standard-Konfiguration):}
\begin{lstlisting}[style=consolestyle]
openssl s_client -connect localhost:4433
\end{lstlisting}

\section{Log-Output Szenario~1}
\label{app:sz1-log}

\begin{lstlisting}[style=consolestyle]
CONNECTED(000001B8)
FC270000:error:0A00042E:SSL routines:ssl3_read_bytes:
tlsv1 alert protocol version:../openssl-3.6.1/ssl/record/
rec_layer_s3.c:918:SSL alert number 70
---
no peer certificate available
\end{lstlisting}

\section{OpenSSL-Befehle Szenario~2}
\label{app:sz2-befehle}

\noindent\textbf{Server-Befehl:}
\begin{lstlisting}[style=consolestyle]
openssl s_server -accept 4433 -cert cert.pem -key key.pem -tls1 \
  -cipher "DEFAULT@SECLEVEL=0"
\end{lstlisting}

\noindent\textbf{Client-Befehl:}
\begin{lstlisting}[style=consolestyle]
openssl s_client -connect localhost:4433 -tls1 \
  -cipher "DEFAULT@SECLEVEL=0"
\end{lstlisting}

\section{Log-Output Szenario~2}
\label{app:sz2-log}

\begin{lstlisting}[style=consolestyle]
CONNECTED(000001BC)
Can't use SSL_get_servername
depth=0 C=AU, ST=Some-State
verify error:num=18:self-signed certificate
verify return:1
---
Certificate chain
 0 s:C=AU, ST=Some-State
   i:C=AU, ST=Some-State
   a:PKEY: RSA, 2048 (bit); sigalg: sha256WithRSAEncryption
---
New, TLSv1.0, Cipher is ECDHE-RSA-AES256-SHA
Protocol : TLSv1
Server public key is 2048 bit
Secure Renegotiation IS NOT supported
\end{lstlisting}


\bibliographystyle{IEEEtran}
\bibliography{IT-Sicherheit_TLS_Literatur}

\end{document}